\section{Resultados}
\label{sec:resultados}

\subsection{Itens 1 e 2: correntes de polarização e de offset}

Com $R_F = 0$, o ponto de operação resultou em $I_{B+}$ =
\SI{0.783}{\nano\ampere} e $v_{out}$ = \SI{-0.7528}{\milli\volt}. Com
$R_F$ = \SI{1}{\mega\ohm}, a medida resultou em $I_{B-}$ =
\SI{1.182}{\nano\ampere} e $v_{out}$ = \SI{0.4293}{\milli\volt}. A
corrente de offset é
\[
  I_{OS} = |I_{B+} - I_{B-}| = |0{,}783 - 1{,}182| = \SI{0.399}{\nano\ampere}.
\]

As duas tensões de saída são consistentes com as expressões previstas
pela análise do circuito quando se inclui o $V_{OS}$ =
\SI{-30.4}{\micro\volt} medido no item 3. Com realimentação negativa e
ganho elevado, a saída se ajusta até que $v_- = v_+ - V_{OS}$ (curto
virtual incluindo o offset). A corrente $I_{B+}$, circulando por $R_I$,
impõe $v_+ = -I_{B+}R_I$; com $R_F = 0$ a saída coincide com $v_-$, e
com $R_F$ = \SI{1}{\mega\ohm} soma-se a queda $I_{B-}R_F$ entre o pino
inversor e a saída. No primeiro caso a saída esperada é
\[
  v_{out} = -I_{B+}R_I - V_{OS}
          = -0{,}783\,\mathrm{mV} + 0{,}030\,\mathrm{mV}
          = -0{,}753\,\mathrm{mV},
\]
igual ao valor simulado. No segundo caso,
\[
  v_{out} = I_{B-}R_F - I_{B+}R_I - V_{OS}
          = 1{,}182 - 0{,}783 + 0{,}030
          = 0{,}429\,\mathrm{mV},
\]
novamente igual ao valor simulado. A folha de dados do OP07C indica
$I_B$ típico de $\pm$\SI{1.8}{\nano\ampere} (máximo
$\pm$\SI{7}{\nano\ampere}) e $I_{OS}$ típico de \SI{0.8}{\nano\ampere}
(máximo \SI{6}{\nano\ampere}), de modo que os valores simulados ficam
dentro das faixas especificadas. Em um amplificador ideal as duas
medidas seriam nulas. O resultado também evidencia a razão do uso de
resistências iguais nas duas entradas: o erro DC é reduzido de $I_B R$
(próximo de \SI{1}{\milli\volt} neste ensaio) para o resíduo $I_{OS}R$,
substancialmente menor.

\subsection{Item 3: tensão de offset de entrada}

A Figura \ref{fig:item3} mostra a curva de $v_{out}$ por $v_{in}$ em
malha aberta. A saída satura próximo de $\pm$\SI{14.2}{\volt} e cruza zero
em $V_{OS}$ = \SI{-30.4}{\micro\volt}. A folha de dados indica típico de
\SI{60}{\micro\volt} e máximo de \SI{150}{\micro\volt} para o OP07C. No
amplificador ideal a curva cruzaria exatamente na origem. A largura da
região linear também é consistente com a teoria: a excursão total da
saída dividida pelo ganho resulta em
\[
  \Delta v_{in} \approx \frac{28{,}4\,\mathrm{V}}{499\,888}
  \approx \SI{57}{\micro\volt},
\]
que é a janela visível na ampliação (direita da Figura \ref{fig:item3}).

\begin{figure}[htb]
  \centering
  \includegraphics[width=0.49\textwidth]{item3_transfer.png}
  \includegraphics[width=0.49\textwidth]{item3_transfer_zoom.png}
  \caption{Item 3: curva de transferência DC em malha aberta, na
  varredura completa (esquerda) e ampliada na região linear (direita).
  O deslocamento do cruzamento por zero é o $V_{OS}$.}
  \label{fig:item3}
\end{figure}

\subsection{Item 4: ganhos DC diferencial e de modo comum}

Com a compensação de \SI{-30.4}{\micro\volt} em série com a entrada, a
curva fica centrada na origem e a derivada numérica $dv_{out}/dv_{in}$
(Figura \ref{fig:item4}) tem valor máximo
\[
  A_{vd} = 499\,888 \approx \SI{500}{\volt\per\milli\volt}
  \qquad (20\log_{10} 499\,888 = \SI{114.0}{\decibel}).
\]
Uma segunda estimativa, obtida por método independente (a inclinação da
reta entre os pontos $v_{in}$ = \SI{-10}{\micro\volt} e
\SI{+10}{\micro\volt}, medida com \texttt{.meas}), resulta no mesmo
valor (499\,888), o que confirma a derivada. O resultado fica
ligeiramente acima do típico do OP07C (\SI{400}{\volt\per\milli\volt},
mínimo \SI{120}{\volt\per\milli\volt}) e coincide com o típico do OP07E
(\SI{500}{\volt\per\milli\volt}).

No caso de modo comum (Figura \ref{fig:item5}), a saída varia pouco em
toda a varredura de $\pm$\SI{3}{\volt} e a derivada resulta em $A_{cm}$
= 0,2494. A inclinação medida entre $\pm$\SI{1}{\volt} com
\texttt{.meas} resulta em 0,2493, valor praticamente idêntico. No
amplificador ideal o ganho diferencial seria infinito e o de modo comum
seria nulo.

\begin{figure}[htb]
  \centering
  \includegraphics[width=0.49\textwidth]{item4_transfer.png}
  \includegraphics[width=0.49\textwidth]{item4_deriv.png}
  \caption{Item 4: curva de transferência com o offset compensado
  (esquerda) e derivada $dv_{out}/dv_{in}$, equivalente ao
  \texttt{dd(u)} do Micro-Cap (direita). O pico da derivada é o
  $A_{vd}$.}
  \label{fig:item4}
\end{figure}

\begin{figure}[htb]
  \centering
  \includegraphics[width=0.49\textwidth]{item5_transfer.png}
  \includegraphics[width=0.49\textwidth]{item5_deriv.png}
  \caption{Item 4, segundo caso: resposta DC ao modo comum (esquerda) e
  derivada $dv_{out}/dv_{cm}$ (direita). O máximo da derivada é o
  $A_{cm}$.}
  \label{fig:item5}
\end{figure}

\subsection{Item 5: CMRR}

Com os dois ganhos medidos, a CMRR é calculada diretamente:
\[
  \mathrm{CMRR} = 20\log_{10}\frac{A_{vd}}{|A_{cm}|}
                = 20\log_{10}\frac{499\,888}{0{,}2494}
                = 20\log_{10}(2{,}00\times 10^{6})
                = \SI{126.0}{\decibel}.
\]
A folha de dados indica CMRR típica de \SI{120}{\decibel} e mínima de
\SI{100}{\decibel} para o OP07C; o ganho diferencial que entra no
cálculo (499\,888) foi comparado à folha de dados na subseção
anterior. No amplificador ideal a CMRR seria infinita. O valor obtido
implica que os $\pm$\SI{3}{\volt} de modo comum do ensaio geram menos
de \SI{3}{\micro\volt} de erro referido à entrada.

\subsection{Item 6: taxa de inflexão (slew rate)}

Na resposta ao degrau (Figura \ref{fig:item6}), a saída parte da
saturação negativa e sobe como uma rampa quase linear até a saturação
positiva, que é o comportamento típico de saída limitada por slew rate.
A medida clássica usa os instantes em que a saída cruza \SI{-10}{\volt}
e \SI{+10}{\volt}, extraídos com \texttt{.meas}:
\[
  \mathrm{SR} = \frac{20\,\mathrm{V}}{t_{hi} - t_{lo}}
              = \frac{20}{105{,}24 - 28{,}41}
              = \frac{20}{76{,}83\,\mu\mathrm{s}}
              = \SI{0.260}{\volt\per\micro\second}.
\]
O patamar da derivada no tempo (mediana na rampa) resulta no mesmo
valor, \SI{0.260}{\volt\per\micro\second}. Ambos são compatíveis com o típico
de \SI{0.3}{\volt\per\micro\second} da folha de dados (mínimo
\SI{0.1}{\volt\per\micro\second}) e também com a estrutura interna do
modelo, em que uma corrente de \SI{7.5}{\micro\ampere} carrega a
capacitância interna de \SI{30}{\pico\farad}: $7{,}5\,\mu\mathrm{A}/30\,\mathrm{pF}$ =
\SI{0.25}{\volt\per\micro\second}. No amplificador ideal a saída
acompanharia o degrau sem atraso.

\begin{figure}[htb]
  \centering
  \includegraphics[width=0.49\textwidth]{item6_step.png}
  \includegraphics[width=0.49\textwidth]{item6_deriv.png}
  \caption{Item 6: degrau de entrada e resposta do seguidor (esquerda) e
  derivada $dv_{out}/dt$, equivalente ao \texttt{ddt(u)} (direita). O
  patamar da derivada é o slew rate.}
  \label{fig:item6}
\end{figure}

\subsection{Item 7 (extra): o mesmo degrau com um amplificador moderno}

Como referência de comparação para o item 6, o mesmo seguidor e o mesmo
degrau foram simulados com um amplificador operacional moderno, o
ADA4610 (entrada JFET, modelo embarcado no LTspice), sem compensação de
offset, dispensável neste ensaio porque o offset desse amplificador é
da ordem de microvolts. A saída cruza de \SI{-10}{\volt} a
\SI{+10}{\volt} em \SI{0,3068}{\micro\second}:
\[
  \mathrm{SR} = \frac{20\,\mathrm{V}}{0{,}3068\,\mu\mathrm{s}}
              = \SI{65.2}{\volt\per\micro\second},
\]
com patamar da derivada em \SI{67.9}{\volt\per\micro\second}. O valor
corresponde a cerca de 250 vezes o slew rate do OP07 no mesmo ensaio
(Figura~\ref{fig:item7}): a transição que no OP07 dura
\SI{77}{\micro\second} se completa em menos de
\SI{0,5}{\micro\second}. A escolha entre os dois componentes depende
dos requisitos de velocidade e de precisão DC da aplicação.

\begin{figure}[htb]
  \centering
  \includegraphics[width=0.49\textwidth]{item6_step.png}
  \includegraphics[width=0.49\textwidth]{item7_ada_step.png}
  \caption{Item 7: o mesmo degrau do item 6 no OP07 (esquerda, escala de
  \SI{200}{\micro\second}) e no ADA4610 (direita, escala de
  \SI{60}{\micro\second}). A rampa do ADA4610 é quase vertical na
  mesma janela.}
  \label{fig:item7}
\end{figure}

\subsection{Verificação das derivadas numéricas}
\label{sec:derivadas}

Como os ganhos e o slew rate são obtidos de derivadas calculadas
numericamente (diferença entre pontos vizinhos dos dados da simulação),
três problemas potenciais foram verificados:

\begin{enumerate}
  \item \textbf{Pico falso no degrau (item 6).} A derivada bruta da
        saída tem um pico de cerca de
        \SI{93000}{\volt\per\micro\second} no instante do degrau de
        entrada. Esse pico não corresponde ao slew rate do
        amplificador: é o salto quase instantâneo da entrada (subida de
        \SI{1}{\pico\second}) que aparece na saída como
        descontinuidade numérica naquele ponto.
        Por isso o cálculo considera apenas a região da rampa, com
        $|v_{out}| < \SI{10}{\volt}$, que é a mesma janela da medida de
        folha de dados.
  \item \textbf{Transitório na saída da saturação (item 6).} Mesmo
        dentro da rampa, o máximo pontual da derivada
        (\SI{0.266}{\volt\per\micro\second}) ocorre logo que o
        amplificador sai da saturação negativa, em $t$ =
        \SI{28.5}{\micro\second}. Trata-se de um transitório curto, não
        do regime da rampa. O valor representativo é o patamar (mediana da
        derivada na rampa), \SI{0.260}{\volt\per\micro\second}, que
        coincide com a medida clássica entre $\pm$\SI{10}{\volt}.
  \item \textbf{Ruído e bordas nas varreduras DC (itens 4 e 5).} Cada
        ponto da varredura provém de uma solução independente do
        simulador, de modo que a derivada ponto a ponto poderia
        amplificar o ruído numérico; além disso o primeiro e o último ponto usam
        diferença unilateral, menos precisa, e por isso não entram na
        procura do máximo. A verificação mostrou que o efeito é pequeno:
        no item 4 o máximo da derivada e a inclinação medida com
        \texttt{.meas} resultam no mesmo 499\,888; no item 5 a derivada
        varia apenas de 0,2492 a
        0,2494 ao longo de toda a varredura (variação de 0,06\%) e a
        inclinação medida resulta em 0,2493. Recalcular a derivada
        usando pontos mais espaçados (5, 25 e 100 passos de distância)
        altera o resultado em menos de 0,05\%.
\end{enumerate}

Cada valor obtido por derivada ($A_{vd}$, $A_{cm}$ e SR) foi, portanto,
confirmado por uma segunda medida independente.

\subsection{Síntese e discussão}

\begin{table}[htb]
  \centering
  \caption{Parâmetros simulados comparados com a folha de dados (OP07C,
  \SI{25}{\celsius}, alimentação de $\pm$\SI{15}{\volt}) e com o
  amplificador ideal.}
  \label{tab:medidas}
  \begin{tabular}{lllll}
    \toprule
    Parâmetro & Simulado & Típico & Máximo & Ideal \\
    \midrule
    $I_{B+}$ & \SI{0.78}{\nano\ampere} & $\pm$\SI{1.8}{\nano\ampere} & $\pm$\SI{7.0}{\nano\ampere} & 0 \\
    $I_{B-}$ & \SI{1.18}{\nano\ampere} & $\pm$\SI{1.8}{\nano\ampere} & $\pm$\SI{7.0}{\nano\ampere} & 0 \\
    $I_{OS}$ & \SI{0.40}{\nano\ampere} & \SI{0.8}{\nano\ampere} & \SI{6.0}{\nano\ampere} & 0 \\
    $V_{OS}$ & \SI{-30.4}{\micro\volt} & \SI{60}{\micro\volt} & \SI{150}{\micro\volt} & 0 \\
    $A_{vd}$ & \SI{500}{\volt\per\milli\volt} & \SI{400}{\volt\per\milli\volt} & sem valor & infinito \\
    $A_{cm}$ & 0,249 & sem valor & sem valor & 0 \\
    CMRR & \SI{126.0}{\decibel} & \SI{120}{\decibel} & sem valor & infinita \\
    SR & \SI{0.260}{\volt\per\micro\second} & \SI{0.3}{\volt\per\micro\second} & sem valor & infinito \\
    \bottomrule
  \end{tabular}
\end{table}

Todos os parâmetros ficaram dentro das faixas do fabricante, o que era
esperado, já que o modelo da Analog Devices é ajustado para
reproduzir os valores típicos da folha de dados (as páginas de
especificações estão no Apêndice \ref{ap:datasheet}). As diferenças restantes têm duas causas principais. Primeiro, o modelo
representa um exemplar fixo do componente, enquanto a folha de dados
descreve a variação estatística da fabricação, razão pela qual ela traz
valores típico e máximo. Segundo, as condições de medida não são
idênticas: o fabricante mede o ganho com carga de \SI{2}{\kilo\ohm} na
saída, e a simulação do item 4 é executada sem carga. Em conjunto, as
medidas justificam a classificação do OP07 como amplificador de
precisão: offset de dezenas de microvolts, correntes de entrada de
poucos nanoamperes e CMRR acima de \SI{100}{\decibel}.
